\documentclass[10pt]{article}

\usepackage[utf8]{inputenc}

\usepackage[T1]{fontenc}

\usepackage{amsmath}

\usepackage{amsfonts}

\usepackage{amssymb}

\usepackage[version=4]{mhchem}

\usepackage{stmaryrd}

\usepackage{CJKutf8}

\usepackage{bbold}

\usepackage{graphicx}

\usepackage[export]{adjustbox}

\graphicspath{ {./images/} }



\begin{document}

форм и прямой суммы изометрий определяют во множестве изометрич. структур операцию структуры $(<,>; T)$ и $\left(<,>^{\prime} ; \frac{T}{T^{\prime}}\right)$ наз. к о б о рд а н т н ы м и, если изометрич. структура $(<,>; T) \perp$ $\perp\left(-<,>^{\prime} ; T^{\prime}\right)$ кобордантна нулю. Пусть $G_{F}=$ множество классов кобордизмов изометрических структур $(<,>; T)$, удовлетворяющих условию $\Delta_{T}(1) \times$ \begin{CJK}{UTF8}{mj}久\end{CJK} $\Delta_{T}(-1) \neq 0$, где $\Delta_{T}(t)$ - характеристич. многочлен изометрии $T$. В изучении групп $G_{+}$и $G_{-}$важную роль играют вложения $f_{+}: G_{+} \rightarrow G_{\mathbb{Q}}$ и $\chi_{-}: G_{-} \rightarrow G_{\mathbb{Q}, \text { к-рые }}$ строятся следующим образом. Каждый класс кобордиз-



\includegraphics[max width=\textwidth, center]{2023_07_09_c32f989b63b88e87ffddg-1}

$A$ - невырожденная $\varepsilon$-матрица, то пусть $B=-A^{-1} A^{\prime}$, $Q=A+A^{\prime}$ и $(<,>; T)$ - изометрич. структура, в к-рой форма $<,>$ задается матрицей $Q$, а изометрия $T$ - матрицей $B$. Это сопоставление корректно определяет гомоморфизмы $\chi_{\varepsilon}$ и $\operatorname{ker} \chi_{\varepsilon}=0$.



Пусть $\alpha=(<,>; T)$ - изометрич. структура на векторном пространстве $V$ и $\lambda \in F[t]$. Пусть $V_{\lambda}$ обозначает $\lambda$-примарную компоненту пространства $V$, т. е. $V_{\lambda}=\operatorname{ker} \lambda(T)^{N}$ для большего $N$. Многочлен $\lambda(t)=t^{k}+$ $+a_{1} t^{k-1}+\ldots+a_{n-1} t+1$ наз. в о з в р а т пы м, если $a_{i}=a_{k-i}$ при всех $i$. Для каждого нешриводимого возвратного многочлена $\lambda \in \mathbb{Q}[t]$ через $\varepsilon_{\lambda}(\alpha)$ обозначается приведенный по $\bmod 2$ показатель, с к-рым $\lambda$ входит в характеристич. многочлен $\Delta_{T}$ изометрии $T$. Для каждого возвратного неприводимого в $\mathbb{R}[t]$ многочлена $\lambda \in \mathbb{R}[t]$ обозначается через $\sigma_{\lambda}(\alpha)$ сигнатура сужения формы $<,>$ на $(V \otimes R) \lambda$. Для каждого простого числа $p$ и возвратного неприводимого в $\mathbb{Q}_{p}[t]$ многочлена $\lambda \in \mathbb{Q}_{p}[t]$ пусть сужение $<,>\lambda$ формы $<,>$ на $\left(V \otimes Q_{p}\right) \lambda$, где $Q_{p}$ - поле $p$-адических чисел. Пусть



\[

\mu_{\lambda}^{p}(\alpha)=(-1,1)^{\frac{r(r+3)}{2}}\left(\operatorname{det}<,>_{\lambda}^{p},-1\right)^{r} S\left(<,>_{\lambda}^{p}\right),

\]



где (,) - символ Гильберта в $\mathbb{Q}_{p}, S$ - символ Хассе, $2 r-$ ранг $<,>p$. Две изометрич. структуры кобордантны тогда и только тогда, когда $\varepsilon_{\lambda}(\alpha)=\varepsilon_{\lambda}(\beta), \sigma_{\lambda}(\alpha)=$ $=\sigma_{\lambda}(\beta), \mu_{\lambda}^{p}(\alpha)=\mu_{\lambda}^{p}(\beta)$ для всех $\lambda$ и $p$, для к-рых эти инварианты определены (см. [3], [4]).



Композиция гомоморфизма Левина, гомоморфизма $\chi$ и функций $\varepsilon \lambda, \sigma \lambda, \mu_{\lambda}^{p}$ сопоставляет каждому нечетномерному узлу $K$ числа $\varepsilon_{\lambda}(K) \in\{0,1\}, \quad \sigma \lambda(K) \in \mathbb{Z}$, $\mu_{\lambda}^{p}(K) \in\{-1,1\}$. Два $(2 q-1)$-мерных узла $K_{1}$ и $K_{2}$, где $q>1$, кобордантны тогда и только тогда, когда $\varepsilon_{\lambda}\left(K_{1}\right)=\varepsilon\left(K_{2}\right), \sigma_{\lambda}\left(K_{1}\right)=\sigma_{\lambda}(K), \mu_{\lambda}^{p}\left(K_{1}\right)=\mu_{\lambda}^{p}\left(K_{2}\right)$ для всех $\lambda$ и $p$, для к-рых эти инварианты определены. $\Sigma \sigma_{\lambda}(K)$ равна сигнатуре узала $K$ (см. Узлов и зацеплений квадратичные бормы), где сумма берется по всем $\lambda(t)$ вида $t^{2}-2 t \cos \theta+1$, где $0<\theta<\pi$, и в этой сумме лишь конечное число слагаемых отлично от нуля.



Аналогичным образом определяются групшы кобордизмов локально плоских и кусочно линейных, к-рые обозначаются $C_{n}^{\text {тор }}$ и $C_{n}^{\mathrm{PL}}$ соответственно. При всех $n$ имеется изоморфизм $C_{n}^{\mathrm{PL}} \approx C_{n}$. Естественное отображение $C_{n} \rightarrow C_{n}^{\text {тор }}$ является изоморфизмом при $n \geqslant 3$, а при $n=3$ оно является мономорфизмом с образом индекса 2. Это, в частности, означает существование несглаживаемых локально плоских топологических трехмерных узлов в $S^{5}$ (см. [5]).



Теория У.к. связана с изучением особенностей не локально плоских и кусочно линейных вложений коразмерности 2. Если $P$ есть $(n+1)$-мерное ориентированное многообразие, вложенное как подкомплекс в $(n+3)$-мерное многообразие $M, x \in P$ и $N$ - малая звездная окрестность $x$ в $M$, то особенность вложения $P$ в $M$ в точке $x$ измеряется следующим образом. Край $\partial N$ является $(n+2)$-мерной сферой, ориентация к-рой определяется ориентацией $M ; P \cap \partial N$ является $n$-мерной сферої, ориентация к-рой определяется ориентацией $P$. Таким образом, возникает $n$-мерный узел $(\partial N, \partial N \cap P)$, к-рый наз. ос об е в нос т ь ю в ло же н и я $P \subset M$ в точке $x$.



Jum.: [1] F o x R. H., M i l n o r J. W., "Osaka J. Math.", France" 1965,93, p $225-71$. [3] L e v in e Jull. Soc. Math. France" "1965, t. 93, p. $425-71 ;$ L L e v i n e J., "Comment.

math. helv.", 1969 , v. 44, p. $229-44 ;$ [4] e r o \% e, "Invent. Math.", 1969, v. 8, p. 98-110; [5] C a p p e l l S , S h a n e s o n

J., "Topology", 1973, v. 12, p 33-40, [6] S t o l t f u s N, “Miem. Amer." Math.'Soc ", 1977, v. 12, №1 192. M. II. Фapбep,



УЗЛОВ ТАБЛИЦА - перечень диаграмм всех простых узлов, допускающих проекции на плоскость с девятью и меньшим тислом двойных точек. Обозначения узлов, приведенных в әтой таблице, стандартны: первая цифра указывает число двойных точек, а вторая (расположенная в пндексе) - порядковый номер узла. Напр., узел $7_{5}$ - әто пятый узел таблицы, имеющий 7 пересечений. Рядом с каждым узлом в закодированном виде указан его многочлен Александера $\Delta(t)=$ $=a_{2 n} t^{2 n}+\ldots+a_{n} t^{n}+\ldots+a_{0}$. Поскольку многочлен Александера всякого узла имеет четную степень и является возвратным (т. е. $a_{i}=a_{2 n-i}$ ), то для его задания достаточно указать набор последних коэфффициентов $a_{n}, a_{n-1}, \ldots, a_{0} ;$ именно они п приведены в таблице. Напр., рядом с узлом 8 у указано $7-5+3-1$. Это означает, что его многочлен Александера равен $\Delta(t)=$ $=-t^{6}+3 t^{5}-5 t^{4}+7 t^{3}-5 t^{2}+3 t-1$. Неальтерни рующие узлы помечены звездочкой. Приводимая таблица (см. колонку 485) перепечатана из [1] с небольшими модификациями.



Jum. [1] B u r d e G, в кH. Jahrbuch Überblicke Mathematik, Mannheim, 1978, S. 131-47. УЗЛОВ ТЕОРИЯ - изучение вложений одномерных многообразий в трехмерное евклидово пространство или в сферу $S^{3}$. В более широком смысле предметом У. т. являются вложения сфер в многообразия (см. Многомерный узел) и вообще вложения многообразиї.



Основные понятия У. т. Вложение (чаще - его образ) несвязной суммы $\mu$ әкземпляров окружности в $\mathbb{R}^{3}$ или $S^{3}$ наз. 3 а це п ле ни е м к р а т нос т и $\mu$. Зацепление кратности $\mu=1$ наз. узл о м. Узлы, составляющие данное зацепление, наз. его к о м п он е н т а м и. Объемлема-изотопич. классы зацеплений (см. Изотопия) наз. т и п а м и зацеплений. Зацепления одного типа наз. э к в и в а ле н т ны м и. Тип зацепления «0, $0, \ldots, 0 » ;$ лежащего в плоскости в $\mathbb{R}^{3}$, наз. т р и в и а л ь н ы м. Несколько простейших узлов имеют специальные названия; напр., узел $3_{1}$ (см. Узлов таблица) наз. т р и л и с т н и к о м, а узел $4_{1}-$ восьмеркой или Листинга узлом. Зацепление, состоящее из нек-рых компонент зацепления $L$, наз. его ч а с т и чны м 3 а ц е п л е н и м. Говорят, что зацешление р а с п а да е т с я (или р а сще пля е т с я), если два его частичных зацепления разделены в $S^{3}$ двумерной сферой. Зацепление наз. 6 р у н н о в ы м, если распадается каждое его частичное зацепление, кроме него самого.



Наиболее изучены кусочно линейные зацепления. Рассмотрение гладких или локально плоских топологич. вложений в $\mathbb{R}^{3}$ приводит к теории, по существу, совпадающей с кусочно линейной.



Обычно зацешления задаются посредством диаграмм (см. Узлов и зацеплений диаграмлы). Если в косе (см. Кос теория) из $2 n$ нитей соединить вверху и внизу по $n$ пар соседних концов отрезками, то получится зацепление, называемое $2 n$-с п л е т е н и е м. Другой способ конструирования зацеплениї из кос состоит в замыкании кос. Если между двумя параллельными плоскостями $\Pi_{1}$ и $\Pi_{2}$ в $\mathbb{R}^{3}$ взять $2 m$ ортогональных им отрезков и соединить их концы попарно $m$ дугами в $\Pi_{1} \cdot$ И $m$ дугами в $\Pi_{2}$ без пересечений, то сумма всех дуг и отрезков даст зацепление. Зацепление, допускающее





\end{document}